\documentclass[letterpaper]{article}
\usepackage{aaai}
\usepackage{times}
\usepackage{helvet}
\usepackage{courier}
\usepackage{amsmath}
\usepackage{amssymb}
\usepackage{graphicx}
\usepackage{booktabs}
\usepackage{algorithm}
\usepackage{algorithmic}
\usepackage{url}
\usepackage{subcaption}
\usepackage{multirow}

\frenchspacing
\setlength{\pdfpagewidth}{8.5in}
\setlength{\pdfpageheight}{11in}

\pdfinfo{
/Title (H-MSB-GRPO: Heuristic Multi-Layer Semantic-Balanced GRPO for Enhanced LLM Reasoning)
/Author (Minh Nguyen Quoc, Bang Nguyen Dinh, Bao Tran Dang)
}

\title{H-MSB-GRPO: Heuristic Multi-Layer Semantic-Balanced GRPO \\ for Enhanced LLM Reasoning}
\author{Minh Nguyen Quoc\textsuperscript{1}, Bang Nguyen Dinh\textsuperscript{1}, Bao Tran Dang\textsuperscript{1}\\
\textsuperscript{1}Ho Chi Minh City University of Technology (HCMUT)\\
District 10, Ho Chi Minh City, Vietnam\\
\{minhnguyen, bang.nguyendinh, bao.trandang\}@hcmut.edu.vn
}

\begin{document}
\maketitle

\begin{abstract}
\begin{quote}
Reinforcement Learning (RL), specifically Group Relative Policy Optimization (GRPO), has emerged as a transformative technique for unlocking deep reasoning in Large Language Models (LLMs). By incentivizing multi-step Chain-of-Thought (CoT) behaviors, GRPO allows models to transcend simple pattern matching. However, current implementations are plagued by critical bottlenecks: advantage vanishing in uniform groups, systematic correction bias during reflection, and inefficient utilization of noisy, high-entropy error samples. We propose \textbf{H-MSB-GRPO} (Heuristic Multi-Layer Semantic-Balanced GRPO), an advanced framework that integrates dynamic multi-layer reflection with semantic error profiling. Our architecture introduces a differentiable \textbf{Heuristic Confidence Gate} that smoothly transitions between standard and negative-enhanced advantage calculation, preventing gradient stalls. We further employ K-Means clustering on reasoning embeddings to identify systematic logical failure modes and construct \textbf{Static Balanced Batches} to stabilize policy gradients. Finally, we derive an information-theoretic loss scaling factor based on the \textbf{Shannon Entropy} of semantic clusters to prioritize fixing systematic errors over random noise. Experimental evaluation on GSM8K and MATH benchmarks shows that H-MSB-GRPO reaches a highly competitive accuracy of \textbf{68.78\%} on the MATH dataset using a 7B model, outperforming vanilla GRPO by \textbf{26.57\%} and its multi-layer predecessors by \textbf{2.46\%}.
\end{quote}
\end{abstract}

\section{Introduction}

The paradigm of Large Language Model (LLM) training has shifted from supervised imitation to reinforcement learning-based alignment. While models like GPT-4 and Claude have shown remarkable reasoning, the release of DeepSeek-R1 popularized Group Relative Policy Optimization (GRPO) \cite{deepseekmath}, which enables specialized reasoning models to be trained without the heavy overhead of a value model. GRPO leverages group-relative rewards, allowing the model to distinguish between "good" and "bad" reasoning paths within the same prompt context.

However, as models scale and tasks become more complex (e.g., Olympiad-level mathematics), vanilla GRPO faces diminishing returns. We identify three fundamental limitations in existing RL pipelines for reasoning:

\textbf{1. The Uniform Batch Problem (Advantage Vanishing).} In many reasoning tasks, especially at the start or late stages of training, a model may either fail all $G$ rollouts in a group or solve them all. In such cases, the mean reward equals the individual reward, and the standard deviation approach zero, causing the advantage signal to vanish. This stalls the gradient and wastes expensive compute.

\textbf{2. The Reflection Paradox (Correction Bias).} Multi-layer or iterative reflection models \cite{mgrpo} often suffer from "over-correction." When prompted to reflect, models often feel "obligated" to change their initial answer, even if it was correct. This bias degrades the model's self-confidence and leads to a high rate of false-negative revisions.

\textbf{3. The Noise vs. Signal Dilemma.} Traditional RL updates treat all errors as equally informative. In reality, an error caused by a simple arithmetic slip (noise) is far less useful than an error caused by a fundamental logical misunderstanding (systematic error). Standard GRPO lacks a mechanism to distinguish between these failure modes.

To address these, we propose \textbf{H-MSB-GRPO} (\textbf{H}euristic \textbf{M}ulti-Layer \textbf{S}emantic-\textbf{B}alanced GRPO). Our framework introduces a two-layer generation-reflection pipeline with a \textbf{Heuristic Confidence Gate}. This gate utilizes a differentiable switching mechanism to apply \textbf{NGRPO} (Negative-enhanced GRPO) \cite{ngrpo} using virtual rewards for uniform batches, while activating a sophisticated \textbf{Semantic Error Profiling (SEP)} module for mixed batches. By clustering reasoning paths and constructing \textbf{Static Balanced Batches} (50/50 correct-error ratio), we ensure the model learns from the most systematic and educative failures.

\section{Related Work}

\subsection{Policy Gradient and GRPO}
Proximal Policy Optimization (PPO) \cite{ppo} has long been the standard for RLHF. However, the requirement for a separate value (critic) model doubles the VRAM usage. GRPO \cite{deepseekmath} solves this by using the average reward of $G$ rollouts as a baseline, effectively simulating a critic without the parameter overhead. NGRPO \cite{ngrpo} later improved this by introducing virtual negative rewards to handle cases where no correct answer is found in the group.

\subsection{Multi-Layer and Iterative Reasoning}
Multi-Layer GRPO (M-GRPO) \cite{mgrpo} proved that self-correction can be learned through a two-pass generation process. By conditioning the second pass on the first, models learn to detect their own inconsistencies. However, M-GRPO's reliance on random selection of training samples often leads to unstable learning curves and correction bias.

\subsection{Semantic Entropy and Balanced Sampling}
Measuring model uncertainty through Semantic Entropy (SE) \cite{kuhn2023semantic} has provided a way to identify when a model is truly confused versus when it is merely being creative. SEED-GRPO \cite{seedgrpo} integrated SE into the loss function to dampen updates for high-uncertainty prompts. Simultaneously, the theory of Balanced Sampling \cite{divagrpo} suggests that training on an equal distribution of correct and "hard" negative samples (the 50/50 ratio) can significantly stabilize RL convergence in sparse reward environments.

\section{Methodology}

\subsection{Formal Derivation of H-MSB-GRPO}
Let $q$ be a prompt and $O = \{o_1, \dots, o_N\}$ be $N$ rollouts from policy $\pi_\theta$. The standard reward for $o_i$ is $r_i \in \{0, 1\}$. In H-MSB-GRPO, we define the accuracy ratio $R = \frac{1}{N} \sum_{i=1}^N r_i$.

\subsubsection{Heuristic Confidence Gate}
To maintain gradient continuity, we define a differentiable heuristic gating function $\Gamma(R)$ using the quadratic form:
\begin{equation}
\Gamma(R) = 4R(1-R)
\end{equation}
Note that $\Gamma \in [0, 1]$, reaching $0$ at $R=0$ (all wrong) or $R=1$ (all correct), and $1$ at $R=0.5$. This parabolic form ensures that purely correct or incorrect batches seamlessly switch to pure NGRPO loss, while balanced batches optimally utilize multi-layer reflection without gradient spikes. This gate selects the advantage calculation method:
\begin{equation}
A_{i} = \Gamma \cdot A_{std} + (1-\Gamma) \cdot A_{ngrpo}
\end{equation}
where $A_{ngrpo}$ uses a virtual reward $r_{virtual} = -1.0$ to inject a negative signal even when all $r_i=0$. The standard advantage is computed as $A_{std} = (r_i - \mu) / \sigma$, where $\mu$ and $\sigma$ are the mean and standard deviation of rewards within the group. When variance approaches zero, $A_{std}$ vanishes, making $A_{ngrpo}$ critical for continuous training.

\subsection{Layered Generation and Reflection}
The pipeline is split into two layers to simulate the "Internal Monologue" and "External Correction":
\begin{enumerate}
    \item \textbf{Layer 1 (Rollout):} $N=8$ samples are generated with temperature $T=0.7$ to explore the solution space.
    \item \textbf{Layer 2 (Reflection):} If $\Gamma > 0$, the model performs $M=2$ reflections per rollout using a neutral system prompt: \textit{"Review the solution logic. Only correct it if errors are found; otherwise, maintain consistency."} This creates a pool of $N \times M = 16$ reasoning paths.
\end{enumerate}

\subsection{Semantic Error Profiling (SEP)}
Incorrect samples are vectorized using a frozen \textit{all-MiniLM-L6-v2} encoder running entirely on the CPU to avoid VRAM contention. We apply K-Means clustering ($K=4$) to partition the errors into semantic clusters $C = \{C_1, \dots, C_k\}$. We identify the \textbf{Top Error Cluster} $C_{top} = \text{argmax}_k |C_k|$. This cluster represents the most common logical failure mode for the given prompt, allowing for \textbf{Targeted Error Correction}.

\subsection{Information-Theoretic Loss Scaling}
We measure the uncertainty of the model's error state via the Shannon Entropy $H$ of the cluster distribution $p_k = |C_k| / \sum |C_j|$:
\begin{equation}
H = - \sum_{k=1}^K p_k \log(p_k)
\end{equation}
The loss scaling factor $\lambda = \exp(-H/\tau)$ (where temperature $\tau=1.0$) is applied to the policy gradient. High entropy (random, diverse arithmetic errors) leads to smaller updates, while low entropy (concentrated, systematic logical errors) leads to larger updates.

\subsection{Hardware-Aware Static Tensor Shaping}
To maximize GPU throughput on NVIDIA H100 and prevent memory fragmentation caused by dynamic sequence lengths, we avoid all dynamic padding. Furthermore, taking inspiration from the DIVA-GRPO variance reduction theorem, we construct a \textbf{Static Balanced Batch} of size $S=16$. We draw $S/2$ samples from the Correct Pool and $S/2$ from $C_{top}$. \textbf{Sampling with Replacement} is used to guarantee the tensor shape is entirely constant at every step. This static configuration is hyper-optimized for maximizing DeepSpeed ZeRO-3 performance.

\begin{algorithm}[H]
\caption{H-MSB-GRPO Training Step}
\begin{algorithmic}[1]
\STATE \textbf{Input:} Prompt $q$, Policy $\pi_\theta$, Verifier $V$
\STATE Generate $N$ rollouts $O = \{o_1, \dots, o_N\}$ via vLLM
\STATE $R \gets \text{accuracy\_ratio}(O, V)$
\STATE $\Gamma \gets 4R(1-R)$
\IF{$\Gamma > 0$}
    \STATE Generate $M$ reflections per rollout $\to$ Pool $P$
    \STATE Cluster $P_{error}$ into $\{C_1 \dots C_K\}$ using K-Means
    \STATE $\lambda \gets \exp(-\text{Entropy}(C)/\tau)$
    \STATE $B \gets$ Sample $S/2$ from $P_{correct}$ and $S/2$ from $C_{top}$
\ELSE
    \STATE $\lambda \gets 1.0$, $B \gets$ Sample $S$ uniformly from $O$
    \STATE Use NGRPO Virtual Advantage $A_{ngrpo}$
\ENDIF
\STATE \textbf{Training Phase:}
\STATE Flush vLLM cache $\to$ Init PyTorch Trainer
\STATE $\rho \gets \frac{\pi_\theta}{\pi_{old}}$
\STATE $\mathcal{L}_{surr} \gets \min\left(\rho A, \text{clip}\left(\rho, 1-\epsilon, 1+\epsilon\right)A\right)$
\STATE $\mathcal{L}_{GRPO} \gets \frac{1}{|B|} \sum \lambda \left[ \mathcal{L}_{surr} - \beta \mathbb{D}_{KL} \right]$
\STATE Backward and Optimizer Step
\end{algorithmic}
\end{algorithm}

\section{Experimental Results}

\subsection{Training Configuration}
We use \textbf{Qwen2.5-7B-Instruct} as the backbone. Training uses \textbf{8x NVIDIA H100} GPUs. We implement a sequential handoff between \textbf{vLLM} (generation) and \textbf{PyTorch/DeepSpeed} (training) to manage the 80GB VRAM constraint. We use \textbf{4-bit QLoRA} with a rank $r=16$.

\subsection{Performance Evaluation}
All models were evaluated using the \texttt{lighteval} framework on the exact same hardware configuration to ensure fairness. Table \ref{tab:main} compares H-MSB-GRPO with vanilla GRPO (which lacks reflection entirely) and M-GRPO (which utilizes reflection but samples revisions uniformly without semantic clustering).

\begin{table}[h]
\centering
\caption{Pass@1 Accuracy (\%) on Reasoning Benchmarks.}
\label{tab:main}
\begin{tabular}{lccc}
\toprule
Dataset & Vanilla GRPO & M-GRPO & \textbf{H-MSB-GRPO} \\
\midrule
GSM8K   & 56.33 & \textbf{70.50} & 69.22 \\
MATH    & 42.21 & 66.32 & \textbf{68.78} \\
\bottomrule
\end{tabular}
\end{table}

\begin{table}[h]
\centering
\caption{MATH Sub-category Analysis: M-GRPO vs H-MSB-GRPO.}
\label{tab:math}
\begin{tabular}{l|cc}
\toprule
Category & M-GRPO (\%) & \textbf{A-MSB (\%)} \\
\midrule
Algebra                     & \textbf{89.38} & 89.13 \\
Counting \& Probability     & 64.55 & \textbf{67.09} \\
Geometry                    & 56.99 & \textbf{57.62} \\
Intermediate Algebra        & 49.83 & \textbf{54.71} \\
Number Theory               & 70.55 & \textbf{73.52} \\
Prealgebra                  & 81.63 & \textbf{82.43} \\
Precalculus                 & 51.28 & \textbf{56.96} \\
\midrule
\textbf{Overall}            & 66.32 & \textbf{68.78} \\
\bottomrule
\end{tabular}
\end{table}

\subsection{Analysis and Case Study}
The results show that H-MSB-GRPO's strength lies in high-complexity reasoning. In the MATH dataset, our method provides a \textbf{2.46\%} absolute gain over M-GRPO. As shown in Table \ref{tab:math}, H-MSB-GRPO yields the most significant improvements in sub-categories requiring deep, multi-step logical derivations, such as \textit{Intermediate Algebra (+4.88\%)} and \textit{Precalculus (+5.68\%)}. For simpler or saturated categories like \textit{Algebra} (~89\%), the semantic profiling dampens unnecessary reflections, maintaining performance parity (-0.25\%). This is largely due to the \textbf{Targeted Error Correction} provided by SEP, proving that K-Means clustering effectively addresses underlying systematic errors rather than just applying pattern matching.

\textbf{Case Study:} For a complex geometry problem, M-GRPO initially generated a correct rollout but "over-corrected" in Layer 2, changing the answer. In contrast, H-MSB-GRPO's \textbf{Heuristic Confidence Gate} and \textbf{Entropy Scaling} identified the high certainty of the Layer 1 response, dampening the Layer 2 update and preserving the correct solution.

\subsection{Ablation Analysis: Deconstructing the Performance Gains}
Because our method builds upon fundamental RL and reflection techniques, our primary baselines inherently serve as functional ablations of the complete H-MSB-GRPO pipeline. By analyzing the performance deltas in Table \ref{tab:main}, we can isolate the contributions of specific modules:

\begin{itemize}
    \item \textbf{Impact of Multi-Layer Reflection (Vanilla GRPO $\rightarrow$ M-GRPO):} Removing the reflection layer entirely (Vanilla GRPO) yields only 42.21\% on the MATH dataset. Introducing a second pass of self-reflection with uniform sampling (M-GRPO) provides a massive +24.11\% boost. This confirms that iterative generation is the core driver of reasoning enhancement. However, sampling revisions blindly limits the model's upper bound due to correction bias.
    \item \textbf{Impact of Semantic Balancing and Guidance (M-GRPO $\rightarrow$ H-MSB-GRPO):} When we replace uniform random sampling in the reflection layer with our proposed Semantic Error Profiling (SEP), Static Balanced Batches, and Entropy Loss Scaling, performance further increases to 68.78\% (+2.46\% over M-GRPO). This isolates the exact contribution of the "Semantic-Balanced" mechanism: intelligently profiling and selecting which errors to learn from is necessary to break through the performance ceiling of naive reflection.
\end{itemize}
This progressive breakdown confirms that while multi-step reflection establishes the basic capacity for reasoning, our dynamic routing and semantic-aware target selection are strictly required to resolve systemic logical bottlenecks efficiently.

\section{Discussion and Limitations}
An intriguing observation from our results is the slight underperformance of H-MSB-GRPO on the GSM8K dataset (69.22\%) compared to M-GRPO (70.50\%). GSM8K consists primarily of simpler arithmetic problems with short reasoning paths. Errors here generally manifest as stochastic calculation slips (arithmetic noise) rather than systemic logical flaws. Consequently, the Semantic Entropy ($H$) for GSM8K errors is uniformly high, causing our Loss Scaling factor ($\lambda$) to excessively dampen gradients. While M-GRPO aggressively updates on these noisy samples to gain a marginal lead on GSM8K, H-MSB-GRPO preserves its representation bounds, leading to vastly superior generalization on the much harder MATH benchmarks.

Additionally, while H-MSB-GRPO is effective, it relies on a frozen embedding model for clustering. If the embedding model fails to capture the nuance of a specific mathematical reasoning step, the SEP module's efficiency degrades. Future work will investigate \textbf{online-learned embeddings} where the vectorizer is co-trained with the policy to better understand domain-specific logical patterns.

\section{Conclusion}
H-MSB-GRPO establishes a robust framework for LLM reasoning by synthesizing multi-layer reflection, semantic error profiling, and hardware-efficient batching. By addressing the fundamental issues of advantage vanishing and correction bias, we achieve state-of-the-art results on challenging mathematical benchmarks.

\bibliographystyle{aaai}
\bibliography{references}

\end{document}

